% !PRINTED 553
\weberpage{553}

Wegen (4) braucht dieser Satz nur bewiesen zu werden für
$\mu_n = n$, weil die Reihe $\Sigma\, \varepsilon_n\theta/n^{s+\frac{1}{2}}$ unbedingt konvergiert und
also nach bekannten elementaren Sätzen diese Eigenschaft hat.

Setzen wir also
\weq{7}{\sigma = \frac{\varepsilon_1}{1^s} + \frac{\varepsilon_2}{2^s} + \frac{\varepsilon_3}{2^s} + \cdots} % [sic] third denominator printed 2^s, not 3^s
und nach (5) $\varepsilon_n = \gamma_n\sqrt{n} - \gamma_{n-1}\sqrt{n-1}$, so wird
\weq{8}{\sigma = \frac{\varepsilon_1}{1^s} + \gamma_2\sqrt{2}\left(\frac{1}{2} - \frac{1}{3^s}\right) + \gamma_3\sqrt{3}\left(\frac{1}{3^s} - \frac{1}{4^s}\right) + \cdots} % [sic] first denominator in the parenthesis is printed $2$, not $2^s$

Da nun für ein unendlich großes $n$
\[
\sqrt{n}\left(\frac{1}{n^s} - \frac{1}{(n+1)^s}\right) = \frac{s}{n^{s+\frac{1}{2}}}
\]
ist, so ist die Reihe (8) unbedingt konvergent, solange $s > \frac{1}{2}$
ist, und daraus folgt auch für diese Reihe unsere Behauptung
nach denselben elementaren Sätzen.

\begin{center}\textbf{\S\,147. Entwickelung von $f, f_1, f_2$.}\end{center}

Die Reihen (12), § 145 sind nun in diesem Falle. Betrachten
wir z.\,B. die zweite von ihnen, aus denen, wie wir gesehen
haben, die anderen hergeleitet werden können, und setzen für
$\mu_1, \mu_2, \mu_3, \ldots$ die der Größe nach geordneten Werte der Funktion
\weq{1}{(x-\alpha y)^2 + \beta^2y^2 = n.}

Wir nehmen $xy$ als rechtwinkelige Koordinaten in der Ebene
und überlagern die Ebene mit zwei \gesperrt{Gittern}, indem wir in dem
einen Gitter für die $x$ die geraden, in dem anderen die ungeraden
ganzen Zahlen setzen.

Die Anzahl der Gitterpunkte, die im Innern der Ellipse (1)
oder auf ihrer Peripherie liegen, für die daher $\mu_n \leqq n$ ist, bezeichnen wir für die beiden Arten mit $Z^0(n), Z^1(n)$.

Setzen wir
\[
x = \xi\sqrt{n}, \qquad y = \eta\sqrt{n},
\]
so geht (1) über in
\weq{2}{(\xi-\alpha\eta)^2 + \beta^2\eta^2 = 1,}
und der Flächeninhalt dieser Ellipse ist $\pi/\beta$, und nach Bd.\,II,
§ 194, 1.:
%CONT>>
% [sic] In(7) steht beim dritten Nenner 2^s statt3^s; so beibehalten.
% [sic] In(8) steht im ersten Klammerterm der Nenner2 statt2^s; so beibehalten.
% [sic] In(8) fehlt beim ersten epsilon-Term der Faktor(1-2^{-s}); gedruckte Formel unveraendert.
